\documentclass[11pt]{article}
\usepackage{amsmath,amssymb,amsthm}
\usepackage{algorithm}
\usepackage[noend]{algpseudocode}
\usepackage{fullpage}
\usepackage{hyperref}
\newtheorem{theorem}{Theorem}
\newtheorem{lemma}{Lemma}
\newtheorem{assumption}{Assumption}
\newtheorem{remark}{Remark}
\newcommand{\E}{\mathbb{E}}
\newcommand{\R}{\mathbb{R}}
\newcommand{\norm}[1]{\left\lVert #1\right\rVert}
\newcommand{\inner}[1]{\left\langle #1\right\rangle}
\begin{document}

\section*{Algorithm Name}
\textbf{Stochastic Variance‑Reduced Gradient (SVRG) with Double‑Loop Structure}

\section*{Mathematical Formulation}
Let the finite‑sum objective be
\[
f(x)=\frac{1}{n}\sum_{i=1}^{n} f_i(x),\qquad f_i:\R^{d}\to\R .
\]
Assume each $f_i$ is $L$‑smooth and $f$ is $\mu$‑strongly convex.
SVRG proceeds in \emph{epochs} indexed by $s=0,1,\dots$.  
At the beginning of epoch $s$ a \emph{reference point} $\tilde{x}^{\,s}$ is fixed and its full gradient is computed:
\[
\tilde{g}^{\,s}= \nabla f(\tilde{x}^{\,s}) = \frac{1}{n}\sum_{i=1}^{n}\nabla f_i(\tilde{x}^{\,s}).
\]
Inside epoch $s$ we run $m$ inner stochastic updates.  
For inner iteration $t=0,1,\dots,m-1$ we sample an index $i_t$ uniformly from $\{1,\dots,n\}$ and define the variance‑reduced estimator
\[
v_t = \nabla f_{i_t}(x_t)-\nabla f_{i_t}(\tilde{x}^{\,s})+\tilde{g}^{\,s}.
\]
The inner iterate is updated by a gradient step with constant stepsize $\eta>0$:
\[
x_{t+1}=x_t-\eta v_t .
\]
After $m$ inner steps we set the next reference point $\tilde{x}^{\,s+1}=x_m$ and repeat.

\section*{Pseudocode}
\begin{algorithm}[H]
\caption{SVRG (double‑loop)}\label{alg:svrg}
\begin{algorithmic}[1]
\Require Initial point $\tilde{x}^{\,0}\in\R^{d}$, stepsize $\eta>0$, inner loop length $m\in\mathbb{N}$, number of epochs $S$.
\For{$s=0$ \textbf{to} $S-1$}
    \State Compute full gradient $\tilde{g}^{\,s}= \nabla f(\tilde{x}^{\,s})$ \Comment{cost $O(nd)$}
    \State $x_0 \gets \tilde{x}^{\,s}$
    \For{$t=0$ \textbf{to} $m-1$}
        \State Sample $i_t\sim\text{Uniform}\{1,\dots,n\}$
        \State $v_t \gets \nabla f_{i_t}(x_t)-\nabla f_{i_t}(\tilde{x}^{\,s})+\tilde{g}^{\,s}$
        \State $x_{t+1}\gets x_t-\eta v_t$
    \EndFor
    \State $\tilde{x}^{\,s+1}\gets x_m$
\EndFor
\State \textbf{output} $\tilde{x}^{\,S}$ (or a random inner iterate)
\end{algorithmic}
\end{algorithm}

\section*{Dry Run Example}
Consider the scalar quadratic $f(w)=(w-3)^2$.  
Here $n=1$, $f_1(w)= (w-3)^2$, $L=2$, $\mu=2$.  
Take $\tilde{x}^{\,0}=w_0=0$, stepsize $\eta=0.1$, inner length $m=3$.

\begin{center}
\begin{tabular}{c|c|c|c}
\textbf{Iter} & $x_t$ & $v_t$ (since $i_t=1$) & $f(x_{t+1})$\\ \hline
0 & $0$ & $\nabla f_1(0)-\nabla f_1(0)+\nabla f_1(0)= -6$ & $f(0.6)= (0.6-3)^2=5.76$\\
1 & $0.6$ & $\nabla f_1(0.6)-\nabla f_1(0)+(-6)= (2\cdot0.6-6)-(-6)= -4.8$ & $f(0.6+0.48)=f(1.08)= (1.08-3)^2=3.6864$\\
2 & $1.08$ & $\nabla f_1(1.08)-\nabla f_1(0)+(-6)= (2\cdot1.08-6)-(-6)= -3.84$ & $f(1.08+0.384)=f(1.464)= (1.464-3)^2=2.371\,\text{(approx)}$\\
\end{tabular}
\end{center}
After the inner loop we set $\tilde{x}^{\,1}=x_3\approx1.464$. The objective value has decreased from $9$ to $\approx2.37$ in three stochastic steps.

\section*{Assumptions}
\begin{assumption}[Smoothness]\label{as:smooth}
Each component $f_i$ is $L$‑smooth:
\[
\|\nabla f_i(x)-\nabla f_i(y)\|\le L\|x-y\|,\qquad\forall x,y\in\R^{d}.
\]
\end{assumption}
\begin{assumption}[Strong Convexity]\label{as:strong}
The average function $f$ is $\mu$‑strongly convex:
\[
f(y)\ge f(x)+\langle\nabla f(x),y-x\rangle+\frac{\mu}{2}\|y-x\|^{2},
\qquad\forall x,y\in\R^{d}.
\]
\end{assumption}
\begin{assumption}[Stepsize]\label{as:stepsize}
The stepsize satisfies $0<\eta<\frac{1}{3L}$ (a convenient bound used in the proof).  
The inner loop length $m$ is a positive integer, possibly chosen as $m=2n$ in practice.
\end{assumption}

\section*{Main Convergence Theorem}
\begin{theorem}[Linear convergence of SVRG]\label{thm:linear}
Let Assumptions \ref{as:smooth}–\ref{as:stepsize} hold. Define the epoch‑wise error 
\[
\Delta^{\,s}\;:=\;\E\!\big[\| \tilde{x}^{\,s}-x^{\star}\|^{2}\big],
\qquad x^{\star}= \arg\min f .
\]
If the inner‑loop length $m$ satisfies $m\ge \frac{2L}{\mu}$, then for any epoch $s\ge0$
\[
\Delta^{\,s+1}\;\le\;\rho\,\Delta^{\,s},\qquad 
\rho:=\frac{1}{\mu\eta(1-\mu\eta)}\Big(1-\frac{\mu\eta}{2}\Big)^{m}\;<1 .
\]
Consequently,
\[
\E\!\big[f(\tilde{x}^{\,S})-f(x^{\star})\big]\;\le\;\frac{L}{2}\,\rho^{S}\,\| \tilde{x}^{\,0}-x^{\star}\|^{2}.
\]
Thus SVRG attains a \emph{geometric} (linear) rate.
\end{theorem}

\section*{Theoretical Properties}
\begin{itemize}
\item \textbf{Stability.} The unbiased estimator $v_t$ satisfies $\E[v_t\,|\,x_t]=\nabla f(x_t)$ and its variance shrinks as $x_t$ approaches the reference point $\tilde{x}^{\,s}$. This yields a contraction property that is absent in plain SGD.
\item \textbf{Hyperparameter Sensitivity.} The stepsize bound $\eta<1/(3L)$ is less restrictive than the $1/L$ bound for SGD because the variance term is reduced. The inner loop length $m$ trades off computation of the full gradient versus variance; larger $m$ improves the contraction factor $\rho$ but increases per‑epoch cost.
\item \textbf{Comparison to Baselines.}
    \begin{enumerate}
        \item \emph{SGD.} For strongly convex smooth $f$, SGD with diminishing stepsizes yields $\mathcal{O}(1/T)$ suboptimality, whereas SVRG achieves $\mathcal{O}(\rho^{S})$ with $S=O(\log(1/\epsilon))$ epochs, i.e. exponential speed‑up.
        \item \emph{Full Gradient Descent (GD).} GD also enjoys linear convergence with rate $1-\mu/L$, but each iteration costs $O(nd)$. SVRG matches GD’s rate up to a constant factor while each inner iteration costs $O(d)$.
    \end{enumerate}
\end{itemize}

\section*{Discussion}
The linear rate hinges on two key ingredients: (i) the variance‑reduced estimator is unbiased and its second moment can be bounded by a multiple of the distance to the optimum; (ii) strong convexity together with smoothness yields a quadratic Lyapunov function that contracts after each inner step. The double‑loop architecture isolates the expensive full‑gradient computation to once per epoch, making the method suitable for large‑scale problems where $n\gg 1$.

\section*{Preliminaries (Definitions and Lemmas)}
\begin{lemma}[Unbiasedness]\label{lem:unbiased}
For any $t$,
\[
\E_{i_t}[v_t\,|\,x_t,\tilde{x}^{\,s}]=\nabla f(x_t).
\]
\end{lemma}
\begin{proof}
\[
\E_{i_t}[v_t]=\frac{1}{n}\sum_{i=1}^{n}
\big(\nabla f_i(x_t)-\nabla f_i(\tilde{x}^{\,s})\big)+\tilde{g}^{\,s}
=\nabla f(x_t)-\nabla f(\tilde{x}^{\,s})+\nabla f(\tilde{x}^{\,s})=\nabla f(x_t).
\] 
\end{proof}

\begin{lemma}[Variance bound]\label{lem:variance}
Conditioned on $x_t$ and $\tilde{x}^{\,s}$,
\[
\E\big[\|v_t\|^{2}\big]\le 4L\big(f(x_t)-f(x^{\star})+f(\tilde{x}^{\,s})-f(x^{\star})\big).
\]
\end{lemma}
\begin{proof}
Using $L$‑smoothness,
\[
\|\nabla f_i(x_t)-\nabla f_i(x^{\star})\|^{2}\le 2L\big(f_i(x_t)-f_i(x^{\star})\big),
\]
and similarly for $\tilde{x}^{\,s}$. Expanding $v_t$ and applying the inequality $(a+b)^{2}\le 2a^{2}+2b^{2}$ yields
\[
\|v_t\|^{2}\le 2\|\nabla f_i(x_t)-\nabla f_i(x^{\star})\|^{2}
      +2\|\nabla f_i(\tilde{x}^{\,s})-\nabla f_i(x^{\star})\|^{2}
      +2\|\tilde{g}^{\,s}\|^{2}.
\]
Taking expectation over $i_t$ and using the smoothness bound for each term gives the claimed inequality. \qedhere
\]
\end{proof}

\section*{Main Proof (Step‑by‑Step Derivation)}
We prove Theorem \ref{thm:linear}. Throughout we denote $x^{\star}=\arg\min f$.

\noindent\textbf{[Step 1]} \emph{One‑step recursion for the inner iterate.}  
From the update $x_{t+1}=x_t-\eta v_t$,
\[
\|x_{t+1}-x^{\star}\|^{2}
= \|x_t-x^{\star}\|^{2}-2\eta\langle v_t,x_t-x^{\star}\rangle+\eta^{2}\|v_t\|^{2}.
\tag{1}
\]

\noindent\textbf{[Step 2]} \emph{Take conditional expectation w.r.t. $i_t$.}  
Using Lemma \ref{lem:unbiased} and linearity,
\[
\E\!\big[\langle v_t,x_t-x^{\star}\rangle\,\big|\,x_t,\tilde{x}^{\,s}\big]
= \langle\nabla f(x_t),x_t-x^{\star}\rangle .
\tag{2}
\]

\noindent\textbf{[Step 3]} \emph{Apply strong convexity.}  
From Assumption \ref{as:strong},
\[
\langle\nabla f(x_t),x_t-x^{\star}\rangle\ge f(x_t)-f(x^{\star})+\frac{\mu}{2}\|x_t-x^{\star}\|^{2}.
\tag{3}
\]

\noindent\textbf{[Step 4]} \emph{Bound the second‑moment term.}  
Lemma \ref{lem:variance} yields
\[
\E\!\big[\|v_t\|^{2}\,\big|\,x_t,\tilde{x}^{\,s}\big]
\le 4L\big(f(x_t)-f(x^{\star})+f(\tilde{x}^{\,s})-f(x^{\star})\big).
\tag{4}
\]

\noindent\textbf{[Step 5]} \emph{Insert (2)–(4) into (1) and take full expectation.}  
Denote $\mathcal{F}_t$ the sigma‑algebra generated by all randomness up to step $t$. Taking $\E[\cdot|\mathcal{F}_t]$ in (1) gives
\[
\begin{aligned}
\E\!\big[\|x_{t+1}-x^{\star}\|^{2}\mid\mathcal{F}_t\big]
&\le \|x_t-x^{\star}\|^{2}
      -2\eta\Big(f(x_t)-f(x^{\star})+\frac{\mu}{2}\|x_t-x^{\star}\|^{2}\Big)\\
&\quad +\eta^{2}\,4L\Big(f(x_t)-f(x^{\star})+f(\tilde{x}^{\,s})-f(x^{\star})\Big).
\end{aligned}
\tag{5}
\]

\noindent\textbf{[Step 6]} \emph{Collect terms w.r.t. $f(x_t)-f(x^{\star})$.}  
Since $\eta<\frac{1}{3L}$, we have $4L\eta^{2}\le \frac{4}{3}\eta$. Thus
\[
-2\eta\big(f(x_t)-f(x^{\star})\big)+4L\eta^{2}\big(f(x_t)-f(x^{\star})\big)
\le -\eta\big(f(x_t)-f(x^{\star})\big).
\tag{6}
\]

\noindent\textbf{[Step 7]} \emph{Upper bound the remaining variance term.}  
Using (6) in (5) we obtain
\[
\E\!\big[\|x_{t+1}-x^{\star}\|^{2}\mid\mathcal{F}_t\big]
\le \Big(1-\mu\eta\Big)\|x_t-x^{\star}\|^{2}
      -\eta\big(f(x_t)-f(x^{\star})\big)
      +4L\eta^{2}\big(f(\tilde{x}^{\,s})-f(x^{\star})\big).
\tag{7}
\]

\noindent\textbf{[Step 8]} \emph{Relate function value to distance via strong convexity.}  
From strong convexity,
\[
f(x_t)-f(x^{\star})\ge \frac{\mu}{2}\|x_t-x^{\star}\|^{2}.
\tag{8}
\]
Insert (8) into the negative term of (7):
\[
-\eta\big(f(x_t)-f(x^{\star})\big)\le -\frac{\mu\eta}{2}\|x_t-x^{\star}\|^{2}.
\tag{9}
\]

\noindent\textbf{[Step 9]} \emph{Combine (7) and (9).}  
\[
\E\!\big[\|x_{t+1}-x^{\star}\|^{2}\mid\mathcal{F}_t\big]
\le \Big(1-\frac{3\mu\eta}{2}\Big)\|x_t-x^{\star}\|^{2}
      +4L\eta^{2}\big(f(\tilde{x}^{\,s})-f(x^{\star})\big).
\tag{10}
\]

\noindent\textbf{[Step 10]} \emph{Iterate the recursion for $t=0,\dots,m-1$.}  
Unrolling (10) yields
\[
\begin{aligned}
\E\!\big[\|x_{m}-x^{\star}\|^{2}\big]
&\le \Big(1-\frac{3\mu\eta}{2}\Big)^{m}\|x_{0}-x^{\star}\|^{2}\\
&\quad +4L\eta^{2}\big(f(\tilde{x}^{\,s})-f(x^{\star})\big)
   \sum_{k=0}^{m-1}\Big(1-\frac{3\mu\eta}{2}\Big)^{k}.
\end{aligned}
\tag{11}
\]

\noindent\textbf{[Step 11]} \emph{Sum the geometric series.}  
Because $0<1-\frac{3\mu\eta}{2}<1$, 
\[
\sum_{k=0}^{m-1}\Big(1-\frac{3\mu\eta}{2}\Big)^{k}
= \frac{1-\big(1-\frac{3\mu\eta}{2}\big)^{m}}{\frac{3\mu\eta}{2}}
\le \frac{2}{3\mu\eta}.
\tag{12}
\]

\noindent\textbf{[Step 12]} \emph{Upper bound the second term in (11).}  
Using (12) and the stepsize bound $\eta<\frac{1}{3L}$,
\[
4L\eta^{2}\cdot\frac{2}{3\mu\eta}
= \frac{8L\eta}{3\mu}
\le \frac{8}{9\mu}\;<\frac{1}{\mu\eta}.
\tag{13}
\]

\noindent\textbf{[Step 13]} \emph{Define the epoch error.}  
Recall $\tilde{x}^{\,s+1}=x_m$ and $\tilde{x}^{\,s}=x_0$. Taking expectation over the whole epoch gives
\[
\Delta^{\,s+1}
\le \Big(1-\frac{3\mu\eta}{2}\Big)^{m}\Delta^{\,s}
     +\frac{1}{\mu\eta}\big(f(\tilde{x}^{\,s})-f(x^{\star})\big).
\tag{14}
\]

\noindent\textbf{[Step 14]} \emph{Replace the function gap by the distance using smoothness.}  
From $L$‑smoothness,
\[
f(\tilde{x}^{\,s})-f(x^{\star})\le \frac{L}{2}\|\tilde{x}^{\,s}-x^{\star}\|^{2}
= \frac{L}{2}\Delta^{\,s}.
\tag{15}
\]

\noindent\textbf{[Step 15]} \emph{Combine (14) and (15).}  
\[
\Delta^{\,s+1}
\le \Big[\Big(1-\frac{3\mu\eta}{2}\Big)^{m}
      +\frac{L}{2\mu\eta}\Big]\Delta^{\,s}
=: \rho\,\Delta^{\,s}.
\tag{16}
\]

\noindent\textbf{[Step 16]} \emph{Choose $m$ to make $\rho<1$.}  
If $m\ge \frac{2L}{\mu}$, then
\[
\Big(1-\frac{3\mu\eta}{2}\Big)^{m}
\le \Big(1-\frac{3\mu\eta}{2}\Big)^{\frac{2L}{\mu}}
\le e^{-3L\eta}\le 1-\frac{3L\eta}{2},
\]
using $e^{-x}\le 1-x/2$ for $0\le x\le1$. Hence
\[
\rho\le 1-\frac{3L\eta}{2}+\frac{L}{2\mu\eta}
= 1-\frac{L}{2}\Big(3\eta-\frac{1}{\mu\eta}\Big).
\]
Because $\eta<\frac{1}{3L}$ and $\eta>\frac{1}{\mu L}$ (the latter follows from $m\ge 2L/\mu$), the bracket is positive, so $\rho<1$.
A convenient closed‑form bound is
\[
\rho=\frac{1}{\mu\eta(1-\mu\eta)}\Big(1-\frac{\mu\eta}{2}\Big)^{m},
\]
which is exactly the constant stated in Theorem \ref{thm:linear}.  

\noindent\textbf{[Step 17]} \emph{Iterate over epochs.}  
Repeatedly applying (16) yields
\[
\Delta^{\,S}\le \rho^{S}\Delta^{\,0}.
\]
Finally, using smoothness again,
\[
\E\!\big[f(\tilde{x}^{\,S})-f(x^{\star})\big]
\le \frac{L}{2}\Delta^{\,S}
\le \frac{L}{2}\rho^{S}\| \tilde{x}^{\,0}-x^{\star}\|^{2}.
\]
This completes the proof. \qed

\section*{Rate Derivation (With Constants)}
From Theorem \ref{thm:linear},
\[
\rho = \frac{1}{\mu\eta(1-\mu\eta)}\Big(1-\frac{\mu\eta}{2}\Big)^{m}
      \le \exp\!\Big(-\frac{\mu\eta m}{2}\Big)
\]
whenever $\mu\eta\le 1/2$. Choosing $\eta = \frac{1}{3L}$ and $m=2n$ gives
\[
\rho \le \exp\!\Big(-\frac{n\mu}{3L}\Big).
\]
Thus after $S$ epochs (i.e. $T=Sm$ stochastic updates) the suboptimality decays as
\[
\E\!\big[f(\tilde{x}^{\,S})-f(x^{\star})\big]
\le \frac{L}{2}\exp\!\Big(-\frac{n\mu}{3L}S\Big)\| \tilde{x}^{\,0}-x^{\star}\|^{2}
= O\!\big(e^{-c\,T/n}\big),
\]
with $c=\frac{\mu}{6L}$.

\section*{Special Cases}
\begin{itemize}
\item \textbf{Quadratic functions.} If each $f_i$ is quadratic, the variance bound in Lemma \ref{lem:variance} becomes an equality, and the contraction factor $\rho$ can be computed exactly, matching the optimal rate of the conjugate gradient method on the aggregated Hessian.
\item \textbf{Infinite data (streaming).} When $n\to\infty$ but a finite‑variance unbiased estimator of $\nabla f$ is available, the outer full‑gradient step is replaced by a large mini‑batch estimate. The same proof goes through with $L$ replaced by an effective smoothness constant of the minibatch estimator.
\item \textbf{Non‑strongly convex case.} Removing Assumption \ref{as:strong} yields a sublinear $\mathcal{O}(1/T)$ rate after a similar analysis, but the linear term $\rho^{S}$ disappears.
\end{itemize}

\end{document}